%============================================================
%  app_manifolds.tex — Banach Manifolds: Reference Summary
%  Tensor Formats in Banach Spaces
%  Antonio Falcó, CEU Universities
%============================================================

\chapter{Banach Manifolds: Reference Summary}
\label{app:manifolds}

This appendix is a reference guide to the differential-geometric notions
used in the monograph. The full exposition --- with complete proofs and
examples --- is in Chapter~\ref{chap:banach}. This appendix records
the precise statements for quick lookup, and adds the material on
connections and curvature that is not needed in the main proofs but
is relevant to the open problems of Chapter~\ref{chap:open}.

Standard references: \cite{Lang1999} and~\cite{Upmeier1985}.

%------------------------------------------------------------
\section{Quick-reference index to Chapter~\ref{chap:banach}}
\label{app:manifolds:index}
%------------------------------------------------------------

\begin{center}
\renewcommand{\arraystretch}{1.35}
\begin{tabular}{@{}p{7cm}l@{}}
\toprule
\textbf{Concept} & \textbf{Location} \\
\midrule
Banach manifold, atlas, AT1--AT3
    & Def.~\ref{def:banach_manifold}, \S\ref{sec:banach:manifolds} \\
Tangent space and tangent map
    & Def.~\ref{def:tangent_space}, \S\ref{sec:banach:manifolds} \\
Immersion, split image condition
    & Def.~\ref{def:immersion}, \S\ref{sec:banach:manifolds} \\
Split (complemented) subspaces
    & Prop.~\ref{prop:split_equiv}, \S\ref{sec:banach:complemented} \\
Finite-dimensional subspaces are split
    & Prop.~\ref{prop:finitedim_split}, \S\ref{sec:banach:complemented} \\
Nilpotent algebra $\cL_{(U,W)}(X)$, $\exp(L)=\Id+L$
    & Prop.~\ref{prop:nilpotent_algebra}, \S\ref{sec:banach:complemented} \\
Grassmann--Banach manifold, chart map $\Psi_{U\oplus W}$
    & Thm.~\ref{thm:grassmann_manifold}, \S\ref{sec:banach:grassmann} \\
Tangent space $T_U\Grass{r}{X}\cong\cL(U,W)$
    & Prop.~\ref{prop:grassmann_tangent}, \S\ref{sec:banach:grassmann} \\
Continuous multilinear maps are $\mathcal{C}^\infty$
    & Prop.~\ref{prop:multilinear_smooth}, \S\ref{sec:banach:multilinear} \\
$\mathcal{C}^1 \Leftrightarrow$ analytic (complex)
    & Prop.~\ref{prop:analytic_c1}, \S\ref{sec:banach:multilinear} \\
Tucker fibre bundle, trivialisations, structure group $G_\fr$
    & \S\ref{sec:tucker_manifold:bundle} \\
Quotient construction fails in $\infty$-dim.\
    & \S\ref{sec:tucker_manifold:bundle}, \S\ref{sec:tucker_manifold:history} \\
\bottomrule
\end{tabular}
\end{center}

%------------------------------------------------------------
\section{ODEs on Banach manifolds}
\label{app:manifolds:ode}
%------------------------------------------------------------

The following theorem is used in Chapter~\ref{chap:DF}
(proof of Theorem~\ref{thm:DF_tucker}) but does not appear
explicitly in Chapter~\ref{chap:banach}.

A \emph{smooth vector field} on a $\mathcal{C}^\infty$-Banach manifold
$\mathbb{M}$ is a $\mathcal{C}^\infty$-map
$F: \mathbb{M} \to T\mathbb{M}$ with $F(m) \in T_m\mathbb{M}$.
An \emph{integral curve} of $F$ through $m_0$ is a $\mathcal{C}^1$-curve
$\gamma: (-\varepsilon,\varepsilon) \to \mathbb{M}$ with
$\gamma(0)=m_0$ and $\dot\gamma(t)=F(\gamma(t))$.

\begin{theorem}[Existence and uniqueness; {\cite[Theorem~4.1.5]{Lang1999}}]
\label{thm:app:ode}
Let $\mathbb{M}$ be a $\mathcal{C}^\infty$-Banach manifold and $F$ a
$\mathcal{C}^\infty$ vector field. For every $m_0\in\mathbb{M}$ there
exist $\varepsilon>0$ and a unique integral curve
$\gamma:(-\varepsilon,\varepsilon)\to\mathbb{M}$ through $m_0$.
The flow map $(t,m_0)\mapsto\gamma(t)$ is $\mathcal{C}^\infty$.
\end{theorem}

\begin{remark}
In local coordinates the theorem reduces to Picard--Lindelöf in Banach
spaces (Theorem~\ref{thm:app:picard}). The $\mathcal{C}^\infty$ flow
follows from smooth dependence on initial data; see~\cite[Ch.~4]{Lang1999}.
\end{remark}

%------------------------------------------------------------
\section{The tangent bundle}
\label{app:manifolds:tangentbundle}
%------------------------------------------------------------

The \emph{tangent bundle} of a $\mathcal{C}^p$-manifold $\mathbb{M}$
($p\geq 1$) is
\begin{equation}
\label{eq:app:tangent_bundle}
    T\mathbb{M} \coloneqq \bigsqcup_{m\in\mathbb{M}} T_m\mathbb{M},
\end{equation}
with projection $\pi_{T\mathbb{M}}:(m,v)\mapsto m$ and zero section
$\sigma_0:m\mapsto(m,0)$. Each chart $(M_i,u_i)$ on $\mathbb{M}$
induces a chart $Tu_i:\pi^{-1}(M_i)\to u_i(M_i)\times X_i$,
making $T\mathbb{M}$ a $\mathcal{C}^{p-1}$-Banach manifold modelled
on $X_i\times X_i$.

\begin{proposition}[{\cite[Chapter~3]{Lang1999}}]
\label{prop:app:tangent_bundle}
$T\mathbb{M}$ is a $\mathcal{C}^{p-1}$-Banach manifold; $\pi_{T\mathbb{M}}$
is a $\mathcal{C}^{p-1}$-submersion; $\sigma_0$ is a
$\mathcal{C}^{p-1}$-embedding. For $f:\mathbb{M}\to\mathbb{N}$
of class $\mathcal{C}^p$, the tangent map $Tf:T\mathbb{M}\to T\mathbb{N}$
satisfies $T(g\circ f)=Tg\circ Tf$.
\end{proposition}

\begin{remark}[Tucker manifold]
\label{rem:app:tangent_TM}
The fibre $T_\bv\Tucker{\fr}{\VD}\cong\Emodel{D}$
(Theorem~\ref{thm:Tucker_Banach_Manifold}(i)); the tangent map of
$\mathfrak{i}:\Tucker{\fr}{\VD}\hookrightarrow\VDnorm$ has image
$\ZD{\bv}$ (Proposition~\ref{prop:tangent_space}); and the immersion
theorem (Theorem~\ref{thm:tucker_immersion}) says this map has
closed split image.
\end{remark}

%------------------------------------------------------------
\section{Connections, parallel transport, and curvature}
\label{app:manifolds:connections}
%------------------------------------------------------------

This section is included for completeness and for
Chapter~\ref{chap:open}; none of this material is used in the
proofs of the main text.

\subsection*{Connections on vector bundles}

A \emph{(linear) connection} on $\pi:E\to\mathbb{M}$ is a bilinear
map $\nabla:\mathfrak{X}(\mathbb{M})\times\Gamma(E)\to\Gamma(E)$
satisfying $\nabla_{fX}s=f\nabla_Xs$ and
$\nabla_X(fs)=(Xf)s+f\nabla_Xs$ for all $f\in\mathcal{C}^\infty(\mathbb{M})$.
In a local trivialisation, $\nabla_Xs=X(s)+\omega(X)(s)$ where $\omega$
is the $\cL(F)$-valued connection $1$-form.

\emph{Parallel transport} along $\gamma:[0,1]\to\mathbb{M}$ gives
isomorphisms $P_t:E_{\gamma(0)}\to E_{\gamma(t)}$ by solving
$\nabla_{\dot\gamma}s=0$, $s(0)=e_0$.

\subsection*{Levi-Civita connection and curvature}

If $\mathbb{M}$ is Riemannian (each $T_m\mathbb{M}$ has a smoothly
varying inner product), the unique torsion-free metric connection is
the \emph{Levi-Civita connection}; its curvature tensor is
\begin{equation}
\label{eq:app:curvature_tensor}
    R(X,Y)Z
    = \nabla_X\nabla_YZ - \nabla_Y\nabla_XZ - \nabla_{[X,Y]}Z.
\end{equation}
The \emph{sectional curvature} at $m \in \mathbb{M}$ in the direction
of a two-plane $\Pi = \spann\{X,Y\} \subset T_m\mathbb{M}$ is
\begin{equation}
\label{eq:app:sectional_curvature}
    K(\Pi) = \frac{\langle R(X,Y)Y, X\rangle}
    {\|X\|^2\|Y\|^2 - \langle X,Y\rangle^2}.
\end{equation}

\begin{proposition}[Jacobi fields and curvature,
{\cite[Chapter~10]{Lang1999}}]
\label{prop:app:jacobi}
Let $\gamma$ be a geodesic in a Riemannian Banach manifold $(\mathbb{M},g)$.
A \emph{Jacobi field} $J(t)$ along $\gamma$ satisfies the Jacobi
equation
\begin{equation}
    \nabla_{\dot\gamma}\nabla_{\dot\gamma} J
    + R(\dot\gamma, J)\dot\gamma = 0.
\end{equation}
If $K \leq \kappa$ along $\gamma$, the spread of neighbouring geodesics
is bounded below by the corresponding space-form comparison: for two
geodesics starting at $\gamma(0)$ with angle $\theta$, the distance
after time $t$ is $\geq t\theta \cdot s_\kappa(t)/t$ where $s_\kappa$
is the sine function in the $\kappa$-space-form.
\end{proposition}

\begin{remark}[Sectional curvature and error bounds for Dirac--Frenkel]
\label{rem:app:curvature_DF}
For the Tucker manifold $\Tucker{\fr}{\VD}$ with $\VDnorm =
\bigotimes_j L^2(I_j)$ (Hilbert case), the sectional curvature of the
inherited Riemannian metric from $\VDnorm$ is non-positive in directions
tangent to the orbit of $\Tucker{\fr}{\VD}$ under the action of
$\prod_\alpha U(r_\alpha)$ (unitary group on the minimal subspaces).
Bounded curvature would imply, via Jacobi field estimates, that the
Dirac--Frenkel projected solution cannot deviate from the full solution
faster than a rate controlled by the curvature; this is the geometric
approach to the error bounds of Open Problem~\ref{prob:lind:error_bounds}.
For the density operator manifold $\Drank{r}{H}$, the curvature is
studied in Open Problem~\ref{prob:curvature_entanglement}.
\end{remark}

\subsection*{Holonomy and homotopy groups}

The \emph{holonomy group} $\mathrm{Hol}(m)$ at $m \in \mathbb{M}$ is the
group of parallel transport isomorphisms $P_\gamma: T_m\mathbb{M} \to
T_m\mathbb{M}$ along all loops $\gamma$ based at $m$. By the
Ambrose--Singer theorem, $\mathrm{Hol}(m)$ is a Lie group whose Lie
algebra is generated by the curvature operators $\{R(X,Y) : X, Y \in
T_m\mathbb{M}\}$~\cite{AbrahamMarsdenRatiu1988}.

\begin{proposition}[Holonomy of the Grassmannian,
{\cite[Chapter~3]{MilnorStasheff1974}}]
\label{prop:app:holonomy_grassmann}
For $\Grass{r}{H}$ with $H$ a separable Hilbert space and the Riemannian
metric induced by $\langle \cdot, \cdot\rangle_H$, the holonomy group
at $U \in \Grass{r}{H}$ is $U(r)$ (acting on $T_U\Grass{r}{H} \cong
\cL(U, U^\perp)$ by left multiplication). In particular,
$\Grass{r}{H}$ is an irreducible symmetric space.
\end{proposition}

%------------------------------------------------------------
\section{The cotangent bundle and differential forms}
\label{app:manifolds:cotangent}
%------------------------------------------------------------

The \emph{cotangent bundle} of $\mathbb{M}$ is
\begin{equation}
    T^*\mathbb{M} \coloneqq \bigsqcup_{m \in \mathbb{M}} (T_m\mathbb{M})^*,
\end{equation}
where $(T_m\mathbb{M})^*$ is the Banach dual of the tangent space.
A \emph{$1$-form} on $\mathbb{M}$ is a section $\alpha: \mathbb{M} \to
T^*\mathbb{M}$; in local coordinates, $\alpha = \sum_i \alpha_i(x)\,dx^i$.
The \emph{exterior derivative} $d\alpha$ is a $2$-form defined by
$(d\alpha)(X,Y) = X(\alpha(Y)) - Y(\alpha(X)) - \alpha([X,Y])$.

\begin{remark}[Relevance to the Tucker and Lindblad manifolds]
\label{rem:app:cotangent_lind}
The cotangent bundle $T^*\Tucker{\fr}{\VD}$ is used implicitly when
defining the gradient of a functional $J: \Tucker{\fr}{\VD} \to \RR$
in the non-Hilbert Banach setting: the Fréchet derivative
$dJ(\bv): T_\bv\Tucker{\fr}{\VD} \to \RR$ is a cotangent vector,
and the \emph{Finsler gradient} $\nabla_F J(\bv) \in
T_\bv\Tucker{\fr}{\VD}$ is its pre-image under the duality mapping
$J: T_\bv\Tucker{\fr}{\VD} \to (T_\bv\Tucker{\fr}{\VD})^*$
(Proposition~\ref{prop:app:duality_mapping}). This is the starting
point for Riemannian and Finsler optimisation on tensor manifolds
(Open Problem~\ref{prob:supervised_learning_banach}).

For the density operator manifold $\Drank{r}{H} \subset \cTsa{1}(H)$,
the cotangent space at $\rho$ is a subspace of $(\cTsa{1}(H))^* \cong
\cT{\infty}(H)_{\mathrm{sa}}$ (bounded self-adjoint operators). The
duality pairing $\langle A, \rho\rangle = \Tr(A\rho)$ gives the natural
pairing between observables and states, and the gradient of a functional
$J(\rho) = \Tr(A\rho)$ (linear observable) is the projection of $A$
onto the tangent space of $\Drank{r}{H}$ via $\hat{P}_\rho$.
\end{remark}
